% WARNING: weak mapping between the source problem and the cited textbook sections.
% Triggering fields from lo_mapping: mapping_confidence = LOW; primary_assessment = needs_change;
% supporting_assessment = incomplete; multipart_assessment = multipart_missing_parts.
% Concepts the source problem requires that NO cited section teaches (all are therefore supplied
% on the student worksheet as reference data, framed as the advisor's procedure sheet, with no
% technique named):
%   1. Entering-column rule (most negative entry in the objective row).
%   2. Minimum-ratio test for choosing the pivot row.
%   3. Optimality / stopping criterion (all objective-row entries under decision columns nonnegative).
%   4. Reading a basic feasible solution off a tableau (non-unit columns set to zero).
%   5. Interpretation of slack columns as unused resource quantities.
% The case still exercises only the skills of the source problem; the cited sections ground the
% row-operation mechanics, and the procedure sheet grounds the rest.
\documentclass[11pt]{article}
\usepackage[margin=0.9in]{geometry}
\usepackage{amsmath,amssymb}
\usepackage{enumitem}
\usepackage{tcolorbox}
\usepackage{xcolor}
\usepackage{booktabs}
\usepackage{array}
\usepackage{fancyhdr}
\tcbuselibrary{skins,breakable}

\definecolor{boxbg}{gray}{0.94}
\definecolor{boxline}{gray}{0.55}
\definecolor{keybg}{gray}{0.97}

\pagestyle{fancy}
\fancyhf{}
\renewcommand{\headrulewidth}{0pt}
\fancyfoot[C]{\small\color{gray} Linear Algebra \quad$\cdot$\quad The Copper Kettle Loan Decision \quad$\cdot$\quad Case Study \#1}
\fancyfoot[R]{\small\color{gray}\thepage}

\newtcolorbox{databox}{enhanced, breakable, colback=boxbg, colframe=boxline,
  boxrule=0.6pt, arc=3pt, left=10pt, right=10pt, top=8pt, bottom=8pt}
\newtcolorbox{keybox}{enhanced, breakable, colback=keybg, colframe=boxline,
  boxrule=0.8pt, arc=3pt, left=12pt, right=12pt, top=10pt, bottom=10pt}

\newcommand{\question}[1]{\medskip\noindent\textbf{#1}\par\smallskip}
\newcommand{\hint}[1]{{\small\color{gray}\textit{Hint: #1}}\par}
\newcommand{\skill}[1]{{\small\textit{Skill: #1}}\par}

\begin{document}
{\small NAME: \underline{\hspace{2.6in}}}\par\vspace{6pt}
\begin{center}
{\LARGE\bfseries The Copper Kettle Loan Decision}\\[3pt]
\rule{2.2in}{0.8pt}\\[5pt]
{\itshape a production-and-financing decision \quad$\bullet$\quad Linear Algebra}
\end{center}

\noindent
Nadia Reyes runs Copper Kettle Foods, a two-product bakehouse in Dayton: her classic
stovetop granola and a cherry-almond snack bar, sold through a Saturday farmers-market
stall and a handful of small grocers. She wants a second oven, and Riverline Community
Bank will underwrite the equipment loan only if Copper Kettle can document that its best
achievable weekly operating profit is at least \$1{,}500. Nadia's business advisor set up
the planning grid below to find that best profit, ran out of session time, and left a
one-page procedure sheet for finishing the job. Nadia's cousin Marco has already finished
the update his own way, and he wants the loan application submitted today. Your group
gets the final word.

\begin{databox}
\textbf{The week at Copper Kettle}
\begin{itemize}[leftmargin=*, itemsep=2pt, topsep=3pt]
\item Granola: 2 oven-hours and 2 packing-hours per batch; \$90 profit per batch.
\item Cherry-almond bars: 1 oven-hour and 5 packing-hours per batch; \$40 profit per batch.
\item Available each week: 24 oven-hours and 40 packing-hours.
\item Bank requirement: documented best achievable weekly profit of at least \$1{,}500.
\item In the grid: $g$ and $c$ are batches of granola and bars; $u_1$ and $u_2$ are the
advisor's bookkeeping quantities, recording a plan's \emph{unused} oven-hours and
\emph{unused} packing-hours; $P$ is weekly profit in dollars.
\end{itemize}
\end{databox}

\begin{center}
\textbf{The advisor's planning grid (as she left it)}\\[4pt]
$\begin{array}{l|rrrrr|r}
 & g & c & u_1 & u_2 & P & \text{Week} \\
\hline
\text{Oven row} & 2 & 1 & 1 & 0 & 0 & 24 \\
\text{Packing row} & 2 & 5 & 0 & 1 & 0 & 40 \\
\text{Profit row} & -90 & -40 & 0 & 0 & 1 & 0 \\
\end{array}$
\end{center}

\begin{databox}
\textbf{The advisor's procedure sheet} (treat these rules as given)
\begin{enumerate}[leftmargin=*, itemsep=2pt, topsep=3pt]
\item Find the most negative number in the bottom row. Its column is this round's
\emph{work column}.
\item For each row above the bottom one whose work-column entry is positive, divide the
last-column number by that entry. The row with the \emph{smallest} quotient is the
\emph{work row}.
\item Scale the work row so its work-column entry becomes 1, then add suitable multiples
of the work row to every other row, the bottom one included, until every other entry in
the work column is 0.
\item Reading a plan from a grid: a column holding a single 1 with 0s elsewhere is
\emph{settled}, and its quantity equals the last-column entry of the row containing the 1.
Every quantity whose column is not of that form equals 0.
\item Stopping check: when no bottom-row entry under the four quantity columns
($g$, $c$, $u_1$, $u_2$) is negative, no further round can raise profit, and the bottom
row's last-column entry is the largest weekly profit any plan within the week's hours can
achieve.
\end{enumerate}
\end{databox}

\question{1. Run One Update Round}
Following the procedure sheet exactly, carry out one full update round on the advisor's
grid. State which column you worked on and which rule chose it, show the two quotients
that chose your row, and write out the complete updated grid, bottom row included.
\hint{Rules 1 and 2 make both choices for you; rule 3 is row arithmetic you have done in class.}

\question{2. Read Off the Week's Plan}
Using rule 4, read the plan your updated grid describes: batches of granola, batches of
bars, unused oven-hours, unused packing-hours, and the week's profit. Then say in one
sentence what this plan tells Nadia to do with the bar line.
\hint{Check each column against rule 4 before you write any number down.}

\question{3. Is This as Good as It Gets?}
Apply rule 5 to your updated grid. Could another update round raise the week's profit?
Point to the exact entries that justify your answer, and state the largest weekly profit
any plan within the week's hours can achieve.
\hint{The answer lives entirely in the bottom row; check the sign of every entry under the four quantity columns.}

\question{4. Settle Marco's Number}
At rule 2, Marco chose the packing row instead, ``because 40 hours leaves more room than
24,'' and his finished grid reads a weekly profit of \$1{,}800, comfortably over the
bank's bar. As a group, decide who is right: find the specific number in Marco's grid
that cannot be true of any real production week, and then make the call against the
bank's \$1{,}500 requirement.
\hint{Before trusting any profit figure, check the plan behind it against the week's oven-hours.}

\bigskip
\noindent\textbf{Group recommendation.} Write Nadia 2 to 3 sentences: submit the loan
application this week or hold it, citing the one profit figure that survives your checks
and what would have to change before you would revisit.
\par\vspace{12pt}\noindent\rule{\textwidth}{0.5pt}
\par\vspace{16pt}\noindent\rule{\textwidth}{0.5pt}
\par\vspace{16pt}\noindent\rule{\textwidth}{0.5pt}

\newpage

\section*{Instructor Material}

\subsection*{Analyst's Checklist (optional, for stuck groups)}
\begin{enumerate}[leftmargin=*, itemsep=2pt, topsep=3pt]
\item Circle the most negative bottom-row entry; its column is the work column.
\item Form both quotients (last column divided by work-column entry); the smaller one
picks the work row.
\item Scale the work row first; then clear the work column's other two entries one row at
a time, writing each row operation down before computing it.
\item Before reading any plan, mark every column settled or not settled per rule 4.
\item For the stopping check, scan only the signs in the bottom row under $g$, $c$,
$u_1$, $u_2$.
\item Compare the one surviving profit figure against \$1{,}500 before drafting the
recommendation.
\end{enumerate}

\subsection*{Answer Key}

\begin{keybox}
\textbf{1. Run One Update Round}
\skill{Executing one simplex pivot: entering column by the most-negative objective-row entry, leaving row by the minimum-ratio test, pivot via elementary row operations.}
Work column: $g$, chosen by rule 1 since $-90$ is the most negative bottom-row entry
(compare $-40$). Quotients (rule 2): oven row $24/2 = 12$; packing row $40/2 = 20$;
smallest is 12, so the oven row is the work row.
Row operations: $R_1 \to \tfrac{1}{2}R_1$, then $R_2 \to R_2 - 2R_1$, then
$R_3 \to R_3 + 90R_1$.
\[
\boxed{\begin{array}{l|rrrrr|r}
 & g & c & u_1 & u_2 & P & \text{Week} \\
\hline
\text{Oven row} & 1 & \tfrac{1}{2} & \tfrac{1}{2} & 0 & 0 & 12 \\
\text{Packing row} & 0 & 4 & -1 & 1 & 0 & 16 \\
\text{Profit row} & 0 & 5 & 45 & 0 & 1 & 1080 \\
\end{array}}
\]

\medskip
\textbf{2. Read Off the Week's Plan}
\skill{Reading the basic feasible solution off a tableau: unit columns are basic, all other variables are zero.}
Settled columns: $g$ (unit in the oven row), $u_2$ (unit in the packing row), and $P$.
Columns $c$ and $u_1$ are not settled, so those quantities are 0.
\[
\boxed{\,g = 12, \quad c = 0, \quad u_1 = 0, \quad u_2 = 16, \quad P = \$1{,}080\,}
\]
Interpretation: bake 12 batches of granola and pause the bar line entirely; the oven is
used to the last hour ($u_1 = 0$) while 16 packing-hours sit idle; the week clears
\$1{,}080. Consistency check in scenario terms: $2(12) = 24$ oven-hours (all used) and
$2(12) = 24$ packing-hours ($40 - 24 = 16$ unused).

\medskip
\textbf{3. Is This as Good as It Gets?}
\skill{Applying the simplex optimality criterion to the objective row.}
Bottom-row entries under $g$, $c$, $u_1$, $u_2$ are $0$, $5$, $45$, $0$: none is
negative, so by rule 5 no further round can raise profit.
\[
\boxed{\,P_{\max} = \$1{,}080\,}
\]
No plan within 24 oven-hours and 40 packing-hours earns more.

\medskip
\textbf{4. Settle Marco's Number (trap, resolved)}
\skill{Minimum-ratio rule as a feasibility guard; comparing a valid and an invalid conclusion against the decision threshold.}
Marco pivoted on the packing row (quotient 20) instead of the oven row (quotient 12),
violating rule 2. His row operations, $R_2 \to \tfrac{1}{2}R_2$, $R_1 \to R_1 - 2R_2$,
$R_3 \to R_3 + 90R_2$, produce
\[
\begin{array}{l|rrrrr|r}
 & g & c & u_1 & u_2 & P & \text{Week} \\
\hline
\text{Oven row} & 0 & -4 & 1 & -1 & 0 & -16 \\
\text{Packing row} & 1 & \tfrac{5}{2} & 0 & \tfrac{1}{2} & 0 & 20 \\
\text{Profit row} & 0 & 185 & 0 & 45 & 1 & 1800 \\
\end{array}
\]
The impossible number is $u_1 = -16$: his plan claims \emph{negative sixteen} unused
oven-hours. Equivalently, his plan bakes $g = 20$ batches, which needs $2(20) = 40$
oven-hours when only 24 exist. Rule 2 exists precisely to prevent this: the smaller
quotient (12) is the most the oven can actually support. Both numbers, side by side:
Marco's \$1{,}800 (from an impossible week) versus the honest \$1{,}080. The gap decides
the case: $\$1{,}800 > \$1{,}500$ says apply, but $\$1{,}080 < \$1{,}500$ says wait.
\[
\boxed{\text{Hold the application: } \$1{,}080 < \$1{,}500}
\]

\medskip
\textbf{Verdict.} Do not submit this week. The only defensible figure is \$1{,}080,
which misses the bank's \$1{,}500 requirement by \$420. A closing note groups often find
on their own: the updated grid prices the bottleneck, since profit rises \$45 for each
added oven-hour, so roughly nine and a half more weekly oven-hours would clear the bar.
That is exactly what the second oven would provide, which is why Nadia should strengthen
the business case before reapplying rather than submit an application the numbers cannot
support.
\end{keybox}

\subsection*{Alignment}
\begin{itemize}[leftmargin=*, itemsep=2pt, topsep=3pt]
\item \textbf{PRIMARY:} A First Course in Linear Algebra (Ken Kuttler), Section 1.3,
``Gaussian Elimination.'' Grounds the row-operation mechanics of the update round
(pivot positions, scaling, elimination).
\item \textbf{SUPPORTING:} Understanding Linear Algebra (David Austin), Section 1.4,
``Pivots and their influence on solution spaces.'' Grounds the rule-4 read-off: columns
with pivot positions correspond to determined (basic) quantities, columns without to
quantities set to zero.
\item \textbf{SUPPORTING:} A First Course in Linear Algebra (Ken Kuttler), Section 1.2,
``Systems of Equations, Algebraic Procedures.'' Grounds why the row moves are legitimate
(elementary operations preserve the solution set).
\item \textbf{Source problem learning objective:} Execute one iteration of the simplex
method on a linear-programming tableau: identify the entering variable from the objective
row, choose the pivot row by the minimum-ratio rule, perform the pivot via elementary row
operations, read off the resulting basic feasible solution, and assess optimality from
the bottom-row entries.
\item \textbf{Section objectives exercised:} Kuttler 1.3, ``Use Gaussian Elimination to
solve systems of linear equations'' (the round's row arithmetic); Austin 1.4, ``columns
of the coefficient matrix containing pivot positions correspond to basic variables and
the columns without pivot positions correspond to free variables'' (reading the plan);
Kuttler 1.2, ``Use elementary operations to find the solution to a linear system of
equations'' (validity of the update).
\item \textbf{What makes this case distinct:} the trap is a ratio-rule violation whose
phantom profit (\$1{,}800) and honest maximum (\$1{,}080) straddle the decision threshold
(\$1{,}500), and the wrong grid is refutable in scenario terms (negative unused
oven-hours). Inert distractors: none among the printed numbers; flavor lives only in the
prose (the farmers-market stall, Dayton).
\item \textbf{Weak-mapping note (honest):} lo\_mapping confidence is LOW
(primary\_assessment needs\_change, supporting\_assessment incomplete,
multipart\_assessment multipart\_missing\_parts). No cited section teaches the
simplex-specific rules (entering-column choice, minimum-ratio test, stopping criterion,
plan read-off, slack interpretation); the worksheet supplies all five as reference data
in the advisor's procedure sheet, with no technique named.
\item \textbf{Attribution:} Adapted from ``A First Course in Linear Algebra'' by Ken
Kuttler, used under CC BY 4.0. Adapted from ``Understanding Linear Algebra'' by David
Austin, used under CC BY 4.0.
\end{itemize}

% VERIFICATION
% Q1: entering column: most negative of (-90, -40) is -90 -> g column. Quotients: 24/2 = 12,
%     40/2 = 20 -> oven row. R1/2 = [1, 1/2, 1/2, 0, 0 | 12]. R2 - 2*R1' = [2-2, 5-1, 0-1, 1-0,
%     0, 40-24] = [0, 4, -1, 1, 0 | 16]. R3 + 90*R1' = [-90+90, -40+45, 0+45, 0, 1, 0+1080]
%     = [0, 5, 45, 0, 1 | 1080]. Matches boxed grid.
% Q2: unit columns: g (row 1), u2 (row 2), P (row 3) -> g = 12, u2 = 16, P = 1080; c = 0, u1 = 0.
%     Scenario check: oven 2*12 + 1*0 = 24, u1 = 24 - 24 = 0. Packing 2*12 + 5*0 = 24,
%     u2 = 40 - 24 = 16. Profit 90*12 + 40*0 = 1080. All consistent.
% Q3: bottom row under g, c, u1, u2 = 0, 5, 45, 0, all >= 0 -> stop; Pmax = 1080. Independent
%     vertex check of max 90g + 40c s.t. 2g + c <= 24, 2g + 5c <= 40, g, c >= 0:
%     (0,0) -> 0; (12,0) -> 1080; (0,8) -> 320; intersection 2g+c=24, 2g+5c=40 -> 4c=16, c=4,
%     g=10 -> 900+160 = 1060. Max = 1080 at (12,0). Confirms optimality.
% Q4: Marco pivots on packing row: R2/2 = [1, 5/2, 0, 1/2, 0 | 20]. R1 - 2*R2' =
%     [0, 1-5, 1, 0-1, 0 | 24-40] = [0, -4, 1, -1, 0 | -16]. R3 + 90*R2' =
%     [0, -40+225, 0, 0+45, 1 | 0+1800] = [0, 185, 0, 45, 1 | 1800]. u1 = -16 < 0 infeasible;
%     plan g = 20 needs 40 oven-hours > 24 available. Trap materiality: 1080 < 1500 < 1800,
%     so the wrong and right numbers sit on opposite sides of the threshold; verdict flips.
% Verdict math: 1500 - 1080 = 420; 420 / 45 = 9.33 extra oven-hours (within the range where the
%     $45-per-oven-hour rate holds, since g = (24+d)/2 stays packing-feasible up to d = 16).

% ATTRIBUTION
% Adapted from "A First Course in Linear Algebra" by Ken Kuttler, used under CC BY 4.0.
% Adapted from "Understanding Linear Algebra" by David Austin, used under CC BY 4.0.

% DERIVED PLAYBOOK (review and promote to library if good)
% P1. SCOPE FENCE. In play: one simplex update round executed as elementary row operations
%     (Kuttler 1.3 mechanics, Kuttler 1.2 legitimacy), reading a solution off a grid via unit
%     columns (Austin 1.4 basic-vs-free correspondence), plus the five printed procedure-sheet
%     rules (entering column, minimum ratio, update, read-off, stopping). Banned because they
%     appear in neither the sections nor the printed reference data: duality, two-phase or
%     Big-M setup, degeneracy and cycling, sensitivity analysis as a named method, matrix
%     inverses, determinants, graphical LP solution as a required route.
% P2. GRAIN SIZE. The source problem's four parts map cleanly onto the arc: (i) pivot ->
%     compute (Q1, with part (iv)'s rule-justification folded in, since the choice happens
%     first in practice); (ii) read BFS -> compute + interpret (Q2); (iii) optimality ->
%     interpret (Q3); the added decision layer (Q4 + recommendation) supplies decide. Four
%     questions, each consuming the previous one's output: Q2 and Q3 read Q1's grid; Q4 weighs
%     Q3's maximum against the threshold. 15-20 group minutes: one pivot dominates the budget.
% P3. FULL STRENGTH. Degenerate version: a grid where the pivot element is already 1, no
%     fractions appear, or the objective row needs no update. Full strength used here: pivot
%     element 2 (forces scaling and halves), eliminations in two rows including the objective
%     row, a BFS with a zero-valued decision variable, and an optimality check that must
%     actually be read from signs (entries 0, 5, 45, 0, with the 5 small enough that groups
%     must look rather than assume).
% P4. CLEAN NUMBERS. Convention: all grid entries integers or halves; RHS integers throughout;
%     profit values multiples of $10; quotients in the ratio test integers (12 and 20); final
%     answers integers (12, 16, 1080, 1800). Engineered backward from the target answers
%     P* = 1080 and trap value 1800 straddling the fixed threshold 1500.
% P5. THE TRAP. Candidate misconceptions for exactly these skills: (a) choosing the pivot row
%     by available capacity or largest RHS instead of the minimum ratio; (b) reading the plan
%     by matching variables to rows in order (g = 12, c = 16); (c) declaring the grid optimal
%     or non-optimal by looking at the RHS instead of the bottom-row signs. Chosen: (a), the
%     ratio-rule violation, because it has decision stakes: it yields a clean phantom profit
%     (1800) on the wrong side of the threshold from the honest maximum (1080), and it is
%     refutable inside the scenario (negative unused oven-hours; a 40-oven-hour plan in a
%     24-oven-hour week). The tie in the work column (2 and 2) makes the wrong row choice
%     genuinely tempting.
% P6. AUTHENTIC USERS. Production planners in small food manufacturing (this case); operations
%     analysts scheduling machine time in job shops; agricultural planners allocating acreage
%     between crops; logistics analysts assigning warehouse labor across product lines. Domain
%     input was blank; the problem's own nutrition/servings context was swapped for a small
%     producer to give the decision (a loan gate) real stakes.
% P7. INTERPRETATION DEMAND. Q1's grid update: none demanded beyond correctness (it feeds Q2).
%     Q2: the BFS must be restated as a production plan (pause the bar line, oven exhausted,
%     packing slack idle). Q3: the stopping rule must be restated as "no better week exists
%     within these hours," turning 1080 into a ceiling, not just a number. Q4: the infeasible
%     grid must be translated into scenario impossibility (negative unused oven-hours), and
%     the two profits must be compared against the bank's threshold to produce the verdict.
% P8. REFERENCE DATA. All five lo_mapping missing_concepts are printed in the student data
%     block as the advisor's procedure sheet, technique unnamed: rule 1 = entering-column
%     choice; rule 2 = minimum-ratio test; rule 3 = the pivot as row operations; rule 4 = BFS
%     read-off; rule 5 = optimality criterion; and the data-block line defining u1, u2 as
%     unused hours = slack interpretation. Nothing else is needed: the row arithmetic itself
%     is taught by the cited sections.
\end{document}
